% Appendix B. Machinery moved out of the body as is (loop0009 item 01); what
% the body keeps is decided in item 05. Sources as for the sections it came
% from: paper2/problem_transformations.md, paper2/proof_reductions.md,
% paper2/revised_algorithm.md, paper2/search_soundness.md.
\section{Supporting material}
\label{app:machinery-all}

This appendix holds the material that supports Sections~\ref{sec:equiv}
and~\ref{sec:search} without being needed to follow them.

\subsection{Edge cases and a gap in a published proof}
\label{app:edge}

Three of the published equalities fail on degenerate inputs, and each failure
is proved in Lean. On a nonempty graph with no edges, $\ns=0$ while
$\vs+1=\theta=1$, so both equalities of \citet{LY2002ref9,LY2002ref10} need an
edge~\leanref{edgeless-ns}.
For a matrix with no 1s, $Z=0$ while $\pw+1=1$. Lengauer's Theorem~4 fails at
$K=0$ on an edgeless graph~\leanref{vsg0}.
\citet[Thm~3.9]{LY2002ref6}, $t=\ns$, fails literally on the $1\times1$ matrix
$[1]$~\leanref{tns}.

The published proof of $\vs\le\ns-1$ \citep[Thm~4.1]{LY2002ref10} has a gap. It
orders the vertices by when they first receive a searcher and claims that every
vertex counted at the worst cut still holds a searcher, ``since the vertices in
$D_{i_0}$ are joined to vertices not yet cleared'' (p.~217). A strategy may place
a searcher and remove it again before clearing anything, which recontaminates
nothing. On $K_{1,3}$, placing and removing a searcher on each leaf before
guarding the centre is an optimal recontamination-free strategy for which the
claim fails~\leanref{kp}. The theorem is true. The
repair orders the vertices by clearing time, as our Lean proof does, or uses the
paper's own normal form, its Corollary~2.4.

\subsection{A thirteenth member}
\label{app:thirteenth}

Progressive black-white pebbling, which is not in Table~1, belongs to the exact
core. \citet[Thm~3]{LY2002ref14} gives
$\vs(G)\le K\iff\mathrm{pbw}(G_d)\le K+2$ on the dag $G_d$ built from $G$; his
Theorem~2 gives $\mathrm{pbw}(D)=\vs(D_u)+1$ on a dag $D$, where $D_u$ is the
MOSP graph of the matrix whose column $v$ is $v$ with its predecessors; and the
minimum progressive pebbling number of \citet[Thm~3.1]{LY2002ref10} equals
$\pw+1$. All three are proved in Lean, each with its hypothesis: Lengauer's
Theorem~2 in its instance form fails at $K=1$ on a dag with no arcs, and
$\mathrm{mpb}=\ns$ fails on edgeless graphs~\leanref{pebbling}.

\subsection{Why the published proof of the definite move fails}
\label{app:whyfails}

The proof moves $q$ to the front of a
solution and claims that after each prefix ``we have at most open(q, S) extra
stacks open, but at least close(q, S) extra stacks closed'' (p.~6). At a prefix
$T\supseteq S$ that does not yet contain $q$, moving $q$ forward changes the open
stacks by $b(T\cup X)-b(T) = |N[q]\setminus O(T)|-|X\setminus T|$. At $T=S$ this
is $\open-\close\le0$. But customers of $X$ that the solution has \emph{already}
closed are in $T$ and no longer count as extra stacks closed, while the stacks
$q$ opens may still be new. In the example, a solution may begin with 3, after
which 4 is free: at $T=\{2,3,4\}$, moving 0 forward costs one more stack. A
second, independent slip: even where the conclusion holds, the literal sequence
$S \mathbin{+\!\!+} [q] \mathbin{+\!\!+} R'$ need not be a solution, because a
customer that $q$ frees stays open until its turn; the argument needs the free
moves.

\subsection{The search for smaller counterexamples}
\label{app:minimality}

The smallest graph on which the rule loses the last solution has 14 customers:
there is none on any graph to 9 vertices, by exhaustion, and none on the 88,802
connected graphs on 10 vertices with at most 14 edges. The counterexample was
found by a search over a family of gadgets built from the failed proof, and
minimised by deletion.

\subsection{The filter and the runs are sound}
\label{app:machinery}

\begin{proposition}\label{prop:filter}
At every state with the invariant, for every set $Q$ of refuted old moves, the
filter definite $\to$ subset $\to$ better with both repairs is node-sound~\leanref{filter}.
\end{proposition}
\begin{proof}
If the definite move fires, Theorem~\ref{thm:repaired} applies. Otherwise the
subset rule keeps a survivor with a solution; let $w$ be the earliest. If the
better move dropped $w$, it cited an earlier survivor that the repaired
argument proves solvable, contradicting the choice of $w$.
\end{proof}

\begin{theorem}[runs are sound]\label{thm:runs}
Let the filter be node-sound and return only playable customers. If at the start
of a run the invariant holds, every memo entry is a genuine refutation, and
every old move is refuted, then the run leaves only genuine memo entries, and if
it answers \emph{false} at $S$ then $\neg\Sol_k(S)$. This holds for any child
order and any memo policy, with the old move on or off~\leanref{runs}.
\end{theorem}
The proof is an induction over the run, in the order in which subtrees complete:
every old move a node holds is a sibling subtree that completed with \emph{false}
earlier in the same run, or an inherited one that passed the reinsertion test.
